\documentclass[11pt]{article}
\usepackage[a4paper,margin=27mm]{geometry}
\usepackage{amsmath,amssymb,amsthm}
\usepackage[hidelinks]{hyperref}
\newtheorem{theorem}{Theorem}
\newtheorem{proposition}[theorem]{Proposition}
\newtheorem{lemma}[theorem]{Lemma}
\newcommand{\eps}{\varepsilon}
\title{An exact correction to a derivative-selection lemma\newline
in the Riemann--Siegel computation}
\author{}
\date{20 September 2026\quad Result SZ-20260920-058}
\begin{document}
\maketitle

\section{The printed claim and its actual computational use}
\label{sec:source}
Juan Arias de Reyna's original author source for
arXiv:2201.00342v1 \cite{AriasII}, Section 3.10, label
\texttt{fincreasing}, asserts that
\begin{equation}\label{eq:H}
H(x)=\frac{5^x}{(2\pi)^{(x-1)/2}\Gamma((x+1)/2)}
\left(\frac{\Gamma(x+1)}{\Gamma(x/3+2)}\right)^{1/2}
\end{equation}
is increasing for all $x>0$. Every factor in this expression is
positive on that domain; the square root is the positive square root.
We give an exact counterexample to that assertion as printed.
We then replace its use in the finite derivative-selection step.
The counterexample concerns this computational lemma, not the
Riemann hypothesis. It does not establish that a particular
implementation returns an incorrect zeta value: the global domain
of the lemma is larger than a particular computation's derivative set.

The source's actual task, labels \texttt{problemFp},
\texttt{defeps5} and \texttt{Fmcondition}, is to calculate the
derivatives at the same original real point $p$, for
\begin{equation}\label{eq:indices}
0\le m\le T:=3L-3,\qquad L\ge1\text{ an integer},\qquad |p|\le1.
\end{equation}
The derivative bound $|F^{(m)}(p)|\le F_m$ is explicitly attributed
in Section 3.9 to the author's Part I, formulas (49) and (55)
\cite{AriasI}. The Taylor coefficient estimate recalled in Section
3.11 is attributed to that work's formulas (47) and (56):
\begin{equation}\label{eq:Fclass}
F(z)=\sum_{j=0}^{\infty}c_{2j}z^{2j},\qquad
|c_{2j}|\le\frac{\pi^j}{2^{j+1}j!}.
\end{equation}
These are the source's inputs, not new estimates for a different
function. Part I has not been independently read in this source
intake; its attribution is retained through the fully read Part II.

Keep the original positive $A(\sigma),B_1,a,\eps_4,F_m$ and define
\begin{align}
E_m:=\eps_5(m)
&=\frac{(\pi^2B_1a)^{m/3}\eps_4}{316A(\sigma)}
\left(\frac{m!}{\Gamma(m/3+2)}\right)^{1/2},\label{eq:eps5}\\
\widetilde E_m:=\widetilde\eps_5(m)
&=\min\{4F_m,E_m\}.\label{eq:tilde}
\end{align}
Our outputs approximate $F^{(m)}(p)$ in absolute value with error
strictly less than \eqref{eq:tilde}. No mass, coefficient, point
or error allowance is replaced. In particular \eqref{eq:tilde}
is the minimum in the original \texttt{defeps5}. A later display
in Section 3.11 prints \texttt{max} at the same place; it cannot
equal $E_m$ when $4F_m>E_m$. The minimum resolves that discrepancy
and is the quantity used throughout this proof.

\section{An exact six-step counterexample}\label{sec:counterexample}
\begin{theorem}\label{thm:ratio}
For every real $x>0$,
\begin{equation}\label{eq:ratio}
\left(\frac{H(x+6)}{H(x)}\right)^2
=\frac{46875^2(x+2)(x+4)}
{\pi^6(x+1)(x+3)(x+5)(x+9)}.
\end{equation}
In particular $H(2006)<H(2000)$.
\end{theorem}
\begin{proof}
Euler's positive integral
$\Gamma(s)=\int_0^\infty t^{s-1}e^{-t}\,dt$ for $s>0$
gives $\Gamma(s+1)=s\Gamma(s)$ by integration by parts: both
boundary terms vanish, since $t^se^{-t}\to0$ at infinity and
$t^s\to0$ at zero. Apply that recurrence separately to the
three gamma factors in the square of \eqref{eq:H}:
\begin{align}
\frac{\Gamma(x+7)}{\Gamma(x+1)}
&=\prod_{j=1}^6(x+j),\label{eq:factorialratio}\\
\frac{\Gamma((x+7)/2)}{\Gamma((x+1)/2)}
&=\frac{(x+1)(x+3)(x+5)}8,\label{eq:halfratio}\\
\frac{\Gamma((x+6)/3+2)}{\Gamma(x/3+2)}
&=\frac{(x+6)(x+9)}9.\label{eq:thirdratio}
\end{align}
The factor from the remaining powers is $5^{12}/(2\pi)^6$.
Thus the ratio before cancellation is exactly
\[
\frac{5^{12}}{(2\pi)^6}
\frac{\prod_{j=1}^6(x+j)}
{\bigl((x+1)(x+3)(x+5)/8\bigr)^2
 ((x+6)(x+9)/9)}.
\]
All cancelled factors are positive. The constant is
$5^{12}64\cdot9/64=9\cdot5^{12}=46875^2$;
cancelling one copy of each odd factor and the factor $x+6$
gives \eqref{eq:ratio}.

For an exact bound on $\pi$, the finite geometric identity gives
\[
\frac1{1+t^2}=\sum_{j=0}^7(-1)^jt^{2j}
+\frac{t^{16}}{1+t^2},\qquad 0\le t\le1.
\]
The integral of its final term is strictly positive. Since
$\pi=4\int_0^1(1+t^2)^{-1}\,dt$, direct rational integration gives
\begin{equation}\label{eq:pi}
\pi>4\sum_{j=0}^7\frac{(-1)^j}{2j+1}
=\frac{135904}{45045}>3,
\qquad 135904-3\cdot45045=769>0.
\end{equation}
At the positive integer $x=2000$, equation \eqref{eq:ratio} and
\eqref{eq:pi} now give the strictly rational certificate
\begin{equation}\label{eq:certificate}
\left(\frac{H(2006)}{H(2000)}\right)^2
<\frac{46875^2\cdot2002\cdot2004}
{3^6\cdot2001\cdot2003\cdot2005\cdot2009}
=\frac{9328515625000}{12454264114047}<1.
\end{equation}
The difference of denominator and numerator is exactly
$3125748489047>0$. Positivity of $H$ proves the conclusion.
There are no rounded gamma evaluations in this argument.
\end{proof}

\section{The Taylor result used by the replacement}\label{sec:taylor}
Result SZ-20260920-056 supplies a complete independent proof of
the first technical gamma inequality in \cite{AriasII} and of
its exact Taylor receiver. Its full proof accompanies this note
\cite{GammaProof}%
% BEGIN VERIFIED PUBLIC PROOF LINK
~\href{https://github.com/KokunoYumeto/zeta-function-research-reader/blob/5992ad50c0f2a06921f1fcaded7b4e38097c0739/workbenches/splitzero-tandem/continuations/20260920-original-kernel-mixed-determinants/providers/GAMMA_INEQUALITY.tex\#L1}{[proof]}%
% END VERIFIED PUBLIC PROOF LINK
{}; we state every constant needed here.
For an integer $M\ge1$, take an integer $J\ge12$ with
\begin{align}
q_J:=\frac{(2\pi)^J}{J!}
&\le\frac{\eps_4}{632A(\sigma)},\label{eq:J1}\\
q_J&\le\frac{\eps_4}{632A(\sigma)M}
\left(\frac{\pi^2B_1a\sqrt3e^2}{M^2}\right)^{(M-1)/3}.
\label{eq:J2}
\end{align}
For the original Taylor polynomial
\begin{equation}\label{eq:P}
P_J(z)=\sum_{j=0}^{J-1}c_{2j}z^{2j},
\end{equation}
that proved receiver states
\begin{equation}\label{eq:Taylor}
|F^{(m)}(p)-P_J^{(m)}(p)|<E_m/2,
\qquad 0\le m<M.
\end{equation}
It includes $m=0$ and $M=1$; at zero the intermediate gamma
comparison is equality, not the strict inequality printed in
the source. Both original inequalities \texttt{defJ} remain unchanged.

The elementary tail estimate underlying the receiver is useful
also when every derivative is retained. Absolute local uniform
convergence of \eqref{eq:Fclass} follows from the majorant
$\tfrac12 e^{\pi |z|^2/2}$. Each differentiated series converges
locally uniformly too, since a fixed derivative contributes at
most a polynomial in $j$ to a factorial-decaying majorant.
For $|p|\le1$, the $m$th derivative of $p^{2j}$ has absolute
value at most $m!\binom{2j}{m}\le m!4^j$ when $m\le2j$,
and is zero otherwise. Therefore
\begin{equation}\label{eq:directtail}
|F^{(m)}(p)-P_J^{(m)}(p)|
\le\frac{m!}2\sum_{j=J}^\infty\frac{(2\pi)^j}{j!}
<m!q_J\quad(J\ge12).
\end{equation}
Indeed the ratio of successive summands is at most
$2\pi/(J+1)<1/2$. The bound $\pi<22/7<13/4$ follows by
integrating the positive function
$t^4(1-t)^4/(1+t^2)$ on $[0,1]$; the integral is
$22/7-\pi$. This exact identity is also proved in
\cite{GammaProof}%
% BEGIN VERIFIED PUBLIC PROOF LINK
~\href{https://github.com/KokunoYumeto/zeta-function-research-reader/blob/5992ad50c0f2a06921f1fcaded7b4e38097c0739/workbenches/splitzero-tandem/continuations/20260920-original-kernel-mixed-determinants/providers/GAMMA_INEQUALITY.tex\#L1}{[proof]}%
% END VERIFIED PUBLIC PROOF LINK
{}, Proposition \texttt{prop:taylor}.
Finally $q_{J+1}/q_J=2\pi/(J+1)$, so $q_J\to0$.
Every finite list of strictly positive upper bounds for $q_J$
can thus be met without appealing to \texttt{fincreasing}.

\section{Finite selection with the original minimum}
\label{sec:selection}
\begin{proposition}\label{prop:finite}
In the source's finite derivative set \eqref{eq:indices}, define
\begin{equation}\label{eq:B}
\mathcal B=\{m\in\{0,\ldots,T\}: E_m\le F_m\}.
\end{equation}
If $\mathcal B$ is empty, zero is an admissible output at every
index. Otherwise put $M=1+\max\mathcal B$, choose $J$ by
\eqref{eq:J1}--\eqref{eq:J2}, and set
\begin{equation}\label{eq:output}
U_m=\begin{cases}P_J^{(m)}(p),&m\in\mathcal B,\\
0,&m\notin\mathcal B.
\end{cases}
\end{equation}
For every $0\le m\le T$,
$|F^{(m)}(p)-U_m|<\widetilde E_m$.
For the computed indices the stronger error bound is $E_m/2$.
\end{proposition}
\begin{proof}
If $m\notin\mathcal B$, then $E_m>F_m>0$ and $4F_m>F_m$;
hence $|F^{(m)}(p)|\le F_m<\min(4F_m,E_m)$.
If $m\in\mathcal B$, then $m<M$ and
$E_m\le F_m<4F_m$, so $\widetilde E_m=E_m$ and
\eqref{eq:Taylor} applies. These two disjoint cases exhaust
the finite set, including $T=0$, and prove the result.
There is no requirement that $\mathcal B$ be an initial interval.
In particular a gap between two computed indices causes no
change to the error estimate or to the original $p$.
\end{proof}

The set \eqref{eq:B} is an exact mathematical description.
Its direct use as a floating-point equality test would be
inappropriate. The next version avoids that test entirely.

\begin{proposition}\label{prop:overlap}
At each index use either of the following certified strict tests:
\begin{align}
F_m<E_m&:\quad\text{output zero},\label{eq:zerotest}\\
E_m<2F_m&:\quad\text{compute the Taylor derivative}.
\label{eq:computetest}
\end{align}
At least one test holds for every positive pair $(F_m,E_m)$.
When both hold either choice is valid. Let $\mathcal C$ be
the finite set assigned to computation. If it is nonempty,
choose $M=1+\max\mathcal C$ and $J$ by
\eqref{eq:J1}--\eqref{eq:J2}. The resulting exact polynomial
outputs satisfy the same original allowances
\eqref{eq:tilde}; each computed output has error $<E_m/2$.

More precisely, suppose the numerical input interface supplies
nested rational enclosing intervals for these positive real
parameters, with widths tending to zero. Checking the strict
tests with the interval endpoints terminates after finitely
many refinements at each index, including $E_m=F_m$ and
$E_m=2F_m$.
\end{proposition}
\begin{proof}
If the zero test fails then $E_m\le F_m<2F_m$, so the compute
test holds. At a computed index $E_m<2F_m<4F_m$ gives
$\widetilde E_m=E_m$ and \eqref{eq:Taylor} applies.
At a zero index the proof of Proposition~\ref{prop:finite}
applies unchanged. Empty $\mathcal C$ needs no $M$ and no $J$.

Write enclosing intervals as $[F_m^-,F_m^+]$ and
$[E_m^-,E_m^+]$. A certified zero test is $F_m^+<E_m^-$;
a certified compute test is $E_m^+<2F_m^-$. If a true
strict test has gap $g>0$, convergence of the interval
endpoints makes the associated endpoint gap positive at
some finite refinement. Since one of the true strict
tests holds, at least one endpoint test eventually passes.
There are only $T+1$ indices, so the maximum of their finite
termination stages is finite. Equality of the two positive
parameters does not require an equality decision: at
$E_m=F_m$ the compute gap is $F_m>0$, while at $E_m=2F_m$
the zero gap is $F_m>0$. This statement describes exactly
what an enclosing-interval numerical interface must return;
it does not identify floating approximations as certificates.
\end{proof}

At a computed index the remaining half allowance can be used
for a certified finite evaluation of $P_J^{(m)}(p)$.
For completeness, given a rational enclosing rectangle for
that complex value whose Euclidean diameter is $<E_m/2$,
its rational center differs from the value by $<E_m/2$.
The triangle inequality with \eqref{eq:Taylor} gives total
error $<E_m=\widetilde E_m$. The polynomial has finitely
many terms; complex rational interval additions and
multiplications enclose them, and convergence of the input
enclosures gives convergence of the output rectangle by
continuity of these finite operations. Choosing its width
against a positive rational lower bound for $E_m/2$ makes
this a terminating enclosure refinement whenever those
input enclosures are provided. This is an exact evaluation
interface and error-composition proof, not a certification
of the source's separate coefficient or rounding algorithms.

\subsection{An explicit rational evaluation constant}
\label{sec:evaluation}
The independent derivation accompanying this note supplies
the following finite bound, whose proof is included here.
For the computed index $m$, set
\begin{align}
I_m&=\{j:0\le j<J,\ 2j\ge m\},&
k_{jm}&=2j-m,& d_{jm}&=\frac{(2j)!}{(2j-m)!},\notag\\
U_j&=\frac{(22/7)^j}{2^{j+1}j!},&
K_m&=\sum_{j\in I_m}d_{jm}(1+k_{jm}U_j).
\label{eq:Km}
\end{align}
The rational $U_j$ bounds $|c_{2j}|$ by \eqref{eq:Fclass}
and $\pi<22/7$. If $I_m$ is empty, the derivative of $P_J$
is identically zero and no evaluation is needed. Otherwise
$K_m$ is a strictly positive rational number. Choose a
positive rational $\eta$ with $\eta K_m<E_m/2$ at every
computed nonzero derivative. There are finitely many such
indices and each $E_m>0$, so positive rational lower bounds
for them give an explicit finite choice of $\eta$.

Take rational complex $\widehat c_{2j}$ and a rational
$r\in[-1,1]$ with
$|\widehat c_{2j}-c_{2j}|\le\eta$ and $|r-p|\le\eta$.
A rational approximation to $p$ may be clamped to $[-1,1]$;
this does not increase its distance from any $p\in[-1,1]$.
Coefficient-coordinate errors at most $\eta/2$ give
complex modulus error at most $\eta/\sqrt2<\eta$.
Now form exactly in rational complex arithmetic
\begin{equation}\label{eq:rationaloutput}
y_m=\sum_{j\in I_m}d_{jm}\widehat c_{2j}r^{k_{jm}}.
\end{equation}
For every nonnegative integer $k$ and $r,p\in[-1,1]$,
$|r^k-p^k|\le k|r-p|$: for $k=0$ both terms are one,
and for $k\ge1$ factor the difference into $r-p$ times
a sum of $k$ products of numbers with absolute value at
most one. Hence
\begin{align}
|y_m-P_J^{(m)}(p)|
&\le\sum_{j\in I_m}d_{jm}
\left(|\widehat c_{2j}-c_{2j}|\,|r|^{k_{jm}}
+|c_{2j}|\,|r^{k_{jm}}-p^{k_{jm}}|\right)\notag\\
&\le\eta K_m<E_m/2.\label{eq:evaluationbound}
\end{align}
This proves the required finite evaluation bound with
every degree and coefficient retained. It covers $p=0$,
$p=\pm1$ and every exponent-zero term. Combined with
\eqref{eq:Taylor}, it gives the original total allowance.

The cutoff can likewise be chosen by rational comparisons.
Take positive rational lower bounds for the two right-hand
sides of \eqref{eq:J1}--\eqref{eq:J2}, and let $b>0$ be
no greater than either. Find $J\ge12$ with
$(44/7)^J/J!<b$. This search terminates because the
successive ratio $(44/7)/(J+1)$ eventually is less than
$1/2$. Since $2\pi<44/7$, both original bounds hold.
Computing the first qualifying $J$ is unnecessary; this
strict rational search avoids an equality test at the
cutoff as well as at derivative selection.

\section{A version retaining every derivative}\label{sec:all}
There is also a replacement that performs no derivative
elimination and no comparison of $E_m$ with $F_m$.
Take $M=T+1$ and choose $J\ge12$ satisfying both original
conditions \eqref{eq:J1}--\eqref{eq:J2} and the finite list
\begin{equation}\label{eq:extraJ}
q_J\le\frac{2F_m}{m!},\qquad 0\le m\le T.
\end{equation}
All right-hand sides are positive; $q_J\to0$ proves
existence. Equations \eqref{eq:Taylor} and
\eqref{eq:directtail} give, for every required $m$,
\begin{equation}\label{eq:halfall}
|F^{(m)}(p)-P_J^{(m)}(p)|
<\min\{E_m/2,2F_m\}=\widetilde E_m/2.
\end{equation}
The remaining half is available for the same certified
finite-polynomial evaluation. Thus \eqref{eq:extraJ}
is an explicit terminating choice of polynomial degree,
not an extra mathematical hypothesis about the source
function or an enlarged error allowance. Strict versions
of all these cutoff inequalities can be certified using
convergent enclosing intervals; they eventually hold
because $q_J\to0$ and all targets stay positive.

\section{Scope, proof map, and provenance}\label{sec:scope}
The exact conclusions are the failure of the printed
global-monotonicity claim, the finite selection rule,
the equality-safe overlapping tests, and the full-index
alternative with the original error tolerances. No assertion
about the source's other technical lemma
\texttt{secondinegamma}, a global zero theorem, or the
complete floating-point Riemann--Siegel implementation is
made. The full-domain gamma proof and Taylor bound used
here are retained separately as result056.

Original-source locators in \cite{AriasII} are
Section 3.9, labels \texttt{defeps5}, \texttt{problemFp};
Section 3.10, label \texttt{fincreasing}; and
Section 3.11, labels \texttt{Fmcondition}, \texttt{defJ},
\texttt{PropAprTaylor}. Root read the entire 2968-line
original \texttt{main.tex}, including the bibliography.
Its SHA-256 is
\begin{center}\small\ttfamily
06786e7e7010b4a51bb431d7e885370d243f\\
52956549d535f69391de120d8bdc.
\end{center}
The intact downloaded archive is separately pinned in the
source-reading ledger. No original author file has been
edited and no author macro has been executed.
The reproducible diagram shows the exact tests and maps
in Propositions~\ref{prop:finite} and \ref{prop:overlap};
it does not display invented derivative samples.
The accompanying integer/rational checker verifies finite
certificates and branch arithmetic, not analytic facts
by numerical sampling. All analytic steps used here are
proved above or in the accompanying full result056 proof.

\begin{thebibliography}{9}
\bibitem{AriasII}
Juan Arias de Reyna,
\emph{High Precision Computation of Riemann's Zeta Function
by the Riemann--Siegel Formula, II}, arXiv:2201.00342v1
(2022), original author \LaTeX\ source.
\url{https://arxiv.org/abs/2201.00342v1}.
\bibitem{AriasI}
Juan Arias de Reyna,
\emph{High precision computation of Riemann's zeta function
by the Riemann--Siegel formula, I}, Mathematics of
Computation \textbf{80} (2011), no.~274, 995--1009.
\url{https://doi.org/10.1090/S0025-5718-2010-02426-3}.
Derivative and coefficient estimates are attributed
through Part II, as specified in Section~\ref{sec:source};
an independent reading of Part I is not claimed.
\bibitem{GammaProof}
\emph{A proof of Arias de Reyna's first gamma inequality
and its exact Taylor-cutoff consequence}, programme
result SZ-20260920-056 (2026), complete accompanying
source \texttt{proof.tex}, Theorem \texttt{thm:gamma}
and Proposition \texttt{prop:taylor}. This is the current
independent mathematical derivation, not a human literature
attribution or an independently published prior result.

% BEGIN VERIFIED PUBLIC PROOF LINK
\href{https://github.com/KokunoYumeto/zeta-function-research-reader/blob/5992ad50c0f2a06921f1fcaded7b4e38097c0739/workbenches/splitzero-tandem/continuations/20260920-original-kernel-mixed-determinants/providers/GAMMA_INEQUALITY.tex\#L1}{Source proof}.
% END VERIFIED PUBLIC PROOF LINK
\end{thebibliography}
\end{document}
